%------------------------------------------------------------------------------%
\section{O Número de Euler ($e$)}

Vimos que o gráfico da função exponencial $y = a^x$ sempre intercepta o eixo
das ordenadas no ponto $(0, 1)$, independentemente do valor da base $a$. No
entanto, a inclinação da reta tangente ao gráfico nesse ponto específico muda
conforme alteramos a base.

\emph{Definição Geométrica de $e$}:
Para tornar a definição mais clara, vamos seguir este raciocínio:

Considere a família de funções $f(x) = a^x$.

Se tomarmos $a = 2$, a reta tangente ao gráfico em $(0, 1)$ tem uma
inclinação aproximada de $0,69$.
Se tomarmos $a = 3$, a reta tangente em $(0, 1)$ tem uma inclinação
aproximada de $1,10$.

%------------------------------------------------------------------------------%
\begin{definicao}{Número Euler}{def:numero-euler}
  Existe um único número real, denotado por $e$,
  tal que a inclinação da reta tangente à curva $y = e^x$ no ponto $(0, 1)$
  é exatamente igual a 1.
\end{definicao}
%------------------------------------------------------------------------------%

Este número é irracional e seu valor aproximado é
\[
  e \approx 2,7182818284\dots
\]
A função
\[
    f(x) = e^x
\]
é chamada de
\emph{função exponencial natural}\index{Função!exponencial natural}.
Sua importância reside no fato de que, nesta base específica,
a taxa de crescimento da função é igual ao próprio valor da
função em qualquer ponto.

%------------------------------------------------------------------------------%
\begin{observation}
  Note que, como a inclinação em $(0, 1)$ é 1, a reta tangente nesse ponto é
  simplesmente a reta $y = x + 1$. Para qualquer outra base, essa relação não
  seria tão ``limpa''.
\end{observation}
%------------------------------------------------------------------------------%

%------------------------------------------------------------------------------%